\chapter[Justificativa]{Justificativa}

A informatização de processos operacionais tem se tornado um fator relevante para a competitividade de pequenos negócios. Em lanchonetes, a agilidade no atendimento e a precisão das informações influenciam diretamente a satisfação dos clientes, o desempenho da equipe e a capacidade de controle do gestor. Quando pedidos, produtos e estoque são administrados sem integração, o estabelecimento fica mais exposto a perdas, atrasos, retrabalho e decisões baseadas em percepções incompletas.

Este trabalho se justifica, inicialmente, pela necessidade prática de desenvolver uma solução capaz de apoiar a rotina de uma lanchonete real ou representativa desse tipo de negócio. A proposta busca reduzir a dependência de registros manuais e centralizar informações importantes em um sistema web, permitindo que diferentes perfis de usuário acessem funcionalidades compatíveis com suas responsabilidades. Dessa forma, pretende-se tornar o atendimento mais organizado e oferecer ao gestor melhores condições para acompanhar a operação.

Do ponto de vista acadêmico, o projeto favorece a aplicação integrada de conhecimentos de análise e desenvolvimento de sistemas, banco de dados, engenharia de software, segurança da informação, arquitetura de software e documentação técnica. A construção de uma aplicação com Laravel e MySQL permite explorar conceitos como MVC, ORM, migrações, autenticação, autorização, regras de negócio e testes funcionais. Além disso, o uso de Docker e Git/GitHub aproxima o desenvolvimento das práticas adotadas em ambientes profissionais.

Sob a perspectiva tecnológica, a escolha do Laravel se justifica pela maturidade do ecossistema PHP e pela existência de recursos que reduzem a complexidade de tarefas recorrentes em aplicações web. O MySQL é adequado por ser amplamente utilizado, possuir documentação consolidada e atender às necessidades de armazenamento relacional do sistema. O Docker contribui para evitar diferenças entre ambientes de desenvolvimento, enquanto o Git permite registrar a evolução do código-fonte e apoiar a rastreabilidade das alterações.

Também há uma justificativa social e econômica. Pequenos empreendedores nem sempre dispõem de recursos para adquirir sistemas comerciais robustos ou excessivamente complexos para sua realidade. Uma solução acadêmica, construída a partir de requisitos bem definidos, pode demonstrar caminhos viáveis para modernização de rotinas e melhoria da organização interna. Ainda que o sistema proposto não substitua integralmente plataformas comerciais completas, ele pode atender a necessidades centrais e servir como base para futuras evoluções.

Portanto, a relevância deste pré-projeto está em articular uma demanda operacional concreta com uma proposta técnica fundamentada. O estudo busca produzir um artefato de software útil, documentado e avaliável, contribuindo tanto para a formação acadêmica do autor quanto para a discussão sobre automação de processos em pequenos estabelecimentos alimentícios.
